Este capítulo parte de los trabajos más cercanos a JobAI, que la Tabla~\ref{tab:ea-trabajos} compara, y desde ellos revisa la literatura de cada elemento que esos trabajos combinan y que este trabajo evalúa: la búsqueda léxica que sirve de línea base, RAG y sus componentes, los motores de almacenamiento vectorial, la clasificación de consultas, la evaluación con colecciones de prueba y jueces, la evaluación de sistemas RAG y la latencia percibida. Cada sección termina con su relación con JobAI.

\section{Trabajos relacionados}
\label{sec:estado-arte-tabla}
\label{sec:estado-arte-empleo}
\label{sec:estado-arte-metodo}

Los trabajos que revisa este capítulo provienen de la literatura reunida durante el desarrollo y de una búsqueda estructurada complementaria, hecha el 16 de noviembre de 2026 en Google Scholar, ACM Digital Library, arXiv y Semantic Scholar con nueve cadenas en inglés, una por tema del capítulo, para el periodo 2019--2026; los trabajos anteriores entraron por búsqueda hacia atrás desde los incluidos. De los 1.395 registros únicos, el cribado por título y resumen, que hizo un asistente de IA y verificó el autor en dos rondas a ciegas, dejó 149 candidatos, y cada cita nueva se verificó en Crossref, arXiv o la página oficial de las actas. La búsqueda no pretende ser una revisión sistemática, y la lectura del texto completo fue parcial: se priorizó la línea de empleo, la más cercana a este trabajo, y aun en ella se leyó menos de la mitad de los candidatos. El Anexo~\ref{anx:busqueda} detalla las cadenas, las bases, los criterios y la verificación del cribado.

La Tabla~\ref{tab:ea-trabajos} compara los trabajos de emparejamiento de empleo más cercanos a este, y los párrafos que la siguen los sitúan en sus líneas de investigación. Las cifras son las que reporta cada trabajo y no son comparables entre sí, porque los corpus, las consultas y los juicios difieren.

{\footnotesize
\begin{longtable}{|>{\raggedright\arraybackslash}p{2.2cm}|>{\raggedright\arraybackslash}p{2.6cm}|>{\raggedright\arraybackslash}p{3.3cm}|>{\raggedright\arraybackslash}p{2.9cm}|>{\raggedright\arraybackslash}p{2.6cm}|}
    \caption{Trabajos relacionados en el emparejamiento de empleo}
    \label{tab:ea-trabajos} \\
    \hline
    \textbf{Trabajo} & \textbf{Dirección y corpus} & \textbf{Comparación} & \textbf{Juicios y métricas} & \textbf{Hallazgo} \\
    \hline
    \endfirsthead
    \hline
    \textbf{Trabajo} & \textbf{Dirección y corpus} & \textbf{Comparación} & \textbf{Juicios y métricas} & \textbf{Hallazgo} \\
    \hline
    \endhead
    \citet{Borisyuk2026} & Buscador y reclutador; producción de LinkedIn & Juez LLM destilado en modelos pequeños & Juez LLM de 8B validado con ejemplos humanos ($\kappa$ lineal 0,77); NDCG@10 & El juez LLM sostiene el entrenamiento y la evaluación a escala \\
    \hline
    \citet{Erol2026} & Hoja de vida $\rightarrow$ ofertas; unas 120.000 ofertas de LinkedIn & \textit{Embeddings} frente a reordenamientos y su combinación & Conjunto pequeño anotado a mano; nDCG@10 & El reordenamiento mejora, con potencia estadística limitada \\
    \hline
    \citet{Vyas2026} & Buscador $\rightarrow$ cargos; 1.283 cargos públicos de Nueva York & Combinaciones de BM25, semántica y grafo de habilidades, con latencia & 30 consultas, dos anotadores ($\kappa \approx 0{,}85$); NDCG@10 0,81 & +7\,\% sobre BM25, con intervalos amplios \\
    \hline
    \citet{Hossain2026} & Buscador $\rightarrow$ ofertas, conversacional; 22.410 ofertas, 988 diálogos & BM25 y densa, con y sin filtros duros; inglés y bengalí romanizado & Una oferta relevante anotada por diálogo; Hit@10 & BM25 gana con texto romanizado y la densa en inglés \\
    \hline
    \citet{Srihari2025JobSphere} & Buscador $\rightarrow$ ofertas; portal público de Punyab & RAG multilingüe frente al portal anterior & Sin descripción de los juicios; P@10 0,68 frente a 0,34 & Mejora frente a la búsqueda por palabras clave \\
    \hline
    \citet{Yoruk2025} & Buscador $\rightarrow$ ofertas; 250 participantes & Agente con RAG frente a la búsqueda tradicional & Sin descripción de quién define una oferta adecuada; precisión en las 5 primeras & +42\,\% de precisión en un experimento A/B \\
    \hline
    \citet{Julianto2026} & Reclutador $\rightarrow$ hojas de vida; 247.480 perfiles & BM25, bi-encoder, RAG y \textit{cross-encoder}, por separado y combinados & P@K, NDCG@K, MRR y latencia & El pipeline completo duplica la P@3 del bi-encoder \\
    \hline
    \citet{Fabregat2026} & Oferta $\rightarrow$ hojas de vida; conjunto sintético en inglés y español & Sistemas participantes en TalentCLEF 2026 & Juicios de la tarea; MAP & Panorama de técnicas: etapas múltiples, fusión, reordenamiento \\
    \hline
    \citet{Thuma2026} & Oferta $\rightarrow$ hojas de vida (TalentCLEF) & Densa, BM25 expandido y su fusión & Juicios de la tarea; MAP y MRR & La fusión ayuda en un idioma y perjudica entre idiomas \\
    \hline
    \citet{Tran2026QTPride} & Oferta $\rightarrow$ hojas de vida (TalentCLEF) & Texto completo frente a vistas estructuradas con LLM, fusionadas con RRF & Juicios de la tarea; MAP & Las vistas estructuradas suman poco \\
    \hline
    Este trabajo & Buscador $\rightarrow$ ofertas; 1.068 ofertas reales, inglés y español & Escalera de componentes; BM25 con y sin los mismos filtros & Juez LLM validado con 298 pares humanos; nDCG@10, cumplimiento de restricciones (C@5), fidelidad, latencia & Capítulo~\ref{cap:resultados} \\
    \hline
\end{longtable}
}

La industria del reclutamiento en línea adoptó pronto la búsqueda semántica con representaciones densas. \citet{HaThuc2017} documentan el paso de LinkedIn Talent Search de una búsqueda por palabras clave a una búsqueda por ejemplos: el reclutador indica uno o varios candidatos ideales y el sistema aprende a construir la consulta y el \textit{ranking} a partir de ellos. A mayor escala, \citet{Shen2024} describen cómo LinkedIn incorporó la recuperación aprendida, densa y generativa, a su sistema de búsqueda y recomendación de empleo, y la misma plataforma usa un modelo de lenguaje para interpretar la consulta y extraer de ella atributos estructurados \citep{Liu2025}. Su sistema actual de búsqueda semántica se entrena y se evalúa con un juez LLM de relevancia que aplica la política de producto y que alcanza un $\kappa$ lineal de 0,77 frente a ejemplos etiquetados por personas \citep{Borisyuk2026}; para servirlo en producción, ese juez se destila en modelos pequeños \citep{Behdin2025}. Otros trabajos usan un modelo de lenguaje que genera, a partir del currículo, la descripción de un empleo adecuado para la persona \citep{ZhengGIRL2023}.

Una segunda línea construye representaciones del dominio laboral. \citet{Decorte2021} aumentan un modelo de lenguaje preentrenado con la coocurrencia de habilidades extraídas de vacantes, y su modelo, JobBERT, supera a los codificadores de oraciones genéricos en la normalización de títulos de cargo. \citet{Bocharova2023} siguen la misma idea con VacancySBERT y reportan mejoras de hasta 21,5\,\% sobre modelos genéricos. ESCO, la clasificación multilingüe europea de habilidades, competencias, cualificaciones y ocupaciones, ofrece un vocabulario común para el emparejamiento \citep{leVrang2014}; CareerBERT la usa para llevar currículos y ocupaciones a un mismo espacio de \textit{embeddings} y la evalúa con ofertas reales y con expertos de recursos humanos \citep{Rosenberger2025}.

La evaluación compartida más cercana es la tarea A de TalentCLEF 2026: dada una oferta, ordenar hojas de vida, en inglés y en español, sobre un conjunto sintético preparado por la organización \citep{Fabregat2026}. Varias participaciones combinan, como JobAI, una lista densa y una léxica. La fusión de una búsqueda densa con BM25 expandido mejoró el MAP en inglés y en español, pero empeoró la búsqueda entre idiomas, donde ganó la densa sola \citep{Thuma2026}. Otras fusionaron con RRF varias vistas de la hoja de vida \citep{Poerwanto2026,Tran2026QTPride} o reordenaron una primera etapa densa \citep{Palomino2026}. En QTPride, las vistas estructuradas extraídas con un LLM sumaron poco sobre el texto completo \citep{Tran2026QTPride}.

En la dirección de JobAI, de quien busca hacia las ofertas, las evaluaciones son más escasas y menos controladas. Synapse recupera ofertas con \textit{embeddings} y las reordena; sus autores reportan una mejora del nDCG@10, pero sobre un conjunto de prueba pequeño anotado a mano que, según ellos, limita la potencia estadística \citep{Erol2026}. JobMatchAI combina BM25, búsqueda semántica y un grafo de habilidades, compara cada combinación de canales con BM25 sobre 30 consultas con juicios de dos anotadores y mide su latencia, aunque calcula las métricas antes de aplicar las restricciones duras \citep{Vyas2026}. Un recomendador conversacional en Bangladés compara BM25 y la búsqueda densa con y sin filtros duros, relaja el filtro cuando no deja candidatas, como JobAI, y muestra que BM25 gana con texto romanizado y la densa en inglés; cada diálogo tiene una sola oferta relevante anotada \citep{Hossain2026}. Los sistemas con RAG suelen reportar mejoras frente a la búsqueda por palabras clave sin describir cómo se juzgó la relevancia \citep{Yoruk2025,Srihari2025JobSphere}. En la dirección contraria, la del reclutador, hay ablaciones con BM25 y \textit{cross-encoder} sobre hojas de vida \citep{Julianto2026} y evaluaciones con RAGAS \citep{Tran2025}.

\textbf{Relación con JobAI.} JobAI no entrena representaciones propias: usa un modelo de \textit{embeddings} general y un pipeline RAG con componentes disponibles. Los trabajos más cercanos en diseño son JobMatchAI, por su comparación de canales frente a BM25, y el recomendador de \citet{Hossain2026}, por sus líneas base léxica y densa con los mismos filtros. Entre los trabajos revisados, ninguno mide si las ofertas devueltas cumplen las restricciones de la consulta, y el juez LLM validado contra personas solo aparece en un sistema industrial \citep{Borisyuk2026}. La evaluación de este trabajo reúne esos elementos. Falta, por tanto, una evaluación repetible que mida, en la búsqueda de empleo desde quien busca, cuánto aporta cada componente frente a BM25 y cuánto cuesta, con un juez validado y sus límites declarados. Las secciones siguientes revisan la literatura de cada uno de esos elementos, y el Capítulo~\ref{cap:metodologia} los reúne en un protocolo de evaluación.

\section{Búsqueda léxica y sus límites en el reclutamiento}
\label{sec:estado-arte-keyword}

La búsqueda léxica relaciona una consulta con un documento por los términos que comparten. Su forma más simple exige que el texto de la consulta aparezca literalmente en el documento; su forma de referencia, Okapi BM25 \citep{Robertson2009}, pondera cada término por su rareza en la colección y normaliza por la longitud del documento. Ambas comparten el \emph{desajuste de vocabulario}: un sistema léxico no puede anticipar todas las formas en que dos personas describen el mismo concepto.

En el reclutamiento, el problema tiene dos caras. \citet{HaThuc2017} lo describen desde el reclutador: traducir criterios de contratación a una consulta booleana exige pericia y, aun así, omite candidatos que describen la misma experiencia con otras palabras. \citet{Bocharova2023} lo describen desde el catálogo de vacantes: dos títulos de cargo pueden ser funcionalmente equivalentes sin compartir ningún término, y un sistema de coincidencia exacta los trata como distintos. El resultado son dos errores que quien busca no distingue: ofertas relevantes que no aparecen porque no comparten el vocabulario de la consulta, y ofertas irrelevantes que aparecen por una palabra clave aislada.

Que la búsqueda léxica tenga esta limitación no la convierte en una línea base débil. \citet{Thakur2021} compararon diez arquitecturas de recuperación sobre 18 colecciones de dominios distintos y encontraron que BM25 es una referencia robusta: varios modelos densos no la superaron fuera de su dominio de entrenamiento, y los modelos de reordenamiento obtuvieron los mejores resultados a un alto costo computacional. En documentos financieros, BM25 superó a la recuperación densa del estado del arte \citep{Akarsu2026}. Una evaluación que compare un sistema semántico solo contra una búsqueda por subcadena sobrestima su ventaja.

\textbf{Relación con JobAI.} La plataforma ofrece en su panel una búsqueda por subcadena sobre el título y la empresa, que la evaluación conserva como línea base de la propia plataforma. La comparación principal, sin embargo, se hace contra BM25, también con los mismos filtros que el pipeline (Capítulo~\ref{cap:metodologia}).

\section{RAG}
\label{sec:estado-arte-rag}

La recuperación densa moderna parte de dos trabajos. \citet{Reimers2019} adaptaron BERT con redes siamesas para producir \textit{embeddings} de oraciones comparables con la similitud coseno, lo que redujo de unas 65 horas a unos 5 segundos la búsqueda del par más similar entre 10.000 oraciones. \citet{Karpukhin2020} mostraron que un recuperador con dos codificadores, entrenado con pocas preguntas y pasajes, superaba a BM25 entre 9 y 19 puntos porcentuales en la exactitud de los 20 primeros pasajes en preguntas de dominio abierto.

RAG fue propuesto por \citet{Lewis2020} como una arquitectura que combina un recuperador denso con un modelo generativo de secuencia a secuencia, entrenados de forma conjunta para responder a partir de pasajes recuperados de una memoria no paramétrica, en su caso Wikipedia. El diseño superó a los modelos puramente paramétricos en tareas intensivas en conocimiento, y la separación entre el modelo y el índice permite actualizar el conocimiento del sistema sin reentrenar el modelo.

La revisión de \citet{Gao2023} organiza la evolución posterior en tres niveles. El RAG \emph{ingenuo} indexa, recupera una vez y genera. El RAG \emph{avanzado} agrega optimizaciones antes de la recuperación, como la reformulación de la consulta, y después de ella, como el reordenamiento y la compresión del contexto. El RAG \emph{modular} descompone el pipeline en componentes intercambiables con flujos condicionales, como el enrutamiento entre estrategias de recuperación. La misma revisión trata, entre los recuperadores, la recuperación híbrida que combina representaciones dispersas y densas. Señala como retos abiertos la calidad de la recuperación, porque fragmentos irrelevantes degradan la respuesta aun con un buen generador, y la persistencia de alucinaciones cuando el contexto es insuficiente.

\textbf{Relación con JobAI.} Antes de este trabajo, la recuperación de JobAI era un RAG ingenuo: una sola búsqueda densa sin reordenamiento. Su enrutamiento entre una ruta semántica y una de agregación SQL ya correspondía a un rasgo del RAG modular (Sección~\ref{sec:estado-arte-hibrido}). La Iteración~2 incorpora componentes del RAG avanzado, y la evaluación mide el aporte de cada uno.


\section{Componentes del RAG avanzado}
\label{sec:estado-arte-componentes}

\subsection{Filtros estructurados y búsqueda aproximada con filtros}

Muchas consultas de empleo combinan una necesidad semántica con condiciones exactas: un cargo descrito con palabras libres, con un salario mínimo, en modalidad remota o con un tipo de contrato. Aplicar esas condiciones sobre los $k$ vecinos más cercanos ya recuperados (\emph{post-filtrado}) puede dejar la lista incompleta cuando el filtro es selectivo; aplicarlas antes, sobre toda la colección (\emph{pre-filtrado}), obliga a una búsqueda exacta sobre el subconjunto. Los índices de grafo como HNSW no se diseñaron para este caso. \citet{Gollapudi2023} proponen índices de grafo que incorporan las etiquetas de filtro en su construcción y en su búsqueda, y \citet{Patel2024} proponen ACORN, una extensión de HNSW que admite predicados arbitrarios recorriendo el subgrafo de los puntos que los cumplen. \citet{Sanca2024} estudian las rutas de acceso de la búsqueda mixta vector-relacional y muestran que el plan más eficiente depende de la selectividad del filtro y del volumen de datos.

\textbf{Relación con JobAI.} Las restricciones de la consulta se extraen con el clasificador y se aplican durante la búsqueda, tanto en la densa como en la léxica (Capítulo~\ref{cap:artefacto}). El \textit{benchmark} de almacenes mide cómo responde cada motor a filtros de distinta selectividad.

\subsection{Búsqueda híbrida y fusión de rankings}

La búsqueda híbrida combina una lista léxica y una densa con la intuición de que modelan la relevancia de formas complementarias: la léxica captura coincidencias exactas de términos poco frecuentes, como nombres de tecnologías, y la densa captura equivalencias sin términos comunes. La fusión por rangos recíprocos (RRF) de \citet{Cormack2009} suma, para cada documento, $1/(k + r)$ por cada lista en la que aparece en la posición $r$; no requiere calibrar los puntajes de los dos sistemas ni datos de entrenamiento, y sus autores mostraron que supera a otros métodos de fusión. \citet{Bruch2023}, en cambio, encontraron que RRF es sensible a su parámetro $k$ y que una combinación convexa de los puntajes normalizados la supera dentro y fuera del dominio, con pocos ejemplos para ajustar su único peso. En el emparejamiento de candidatos y ofertas, \citet{Yashwant2026} usan RRF para combinar la similitud vectorial con la búsqueda de texto completo, y en TalentCLEF la fusión densa y léxica ayudó en un mismo idioma y perjudicó entre idiomas \citep{Thuma2026}.

\textbf{Relación con JobAI.} El pipeline fusiona la búsqueda densa y BM25 con RRF y $k = 60$, sin ajustar. Los hallazgos de \citet{Bruch2023} y de \citet{Thuma2026} son una referencia para interpretar el aporte de este componente, en especial en las consultas en español sobre ofertas en inglés.

\subsection{Expansión de la consulta y del documento}

Otra forma de atacar el desajuste de vocabulario es agregar términos a la consulta o al documento. Los modelos de relevancia estiman, a partir de la consulta sola y sin datos de entrenamiento, la probabilidad de cada palabra en la clase relevante \citep{LavrenkoCroft2001}; su variante RM3, que los usa como retroalimentación de relevancia simulada, es la referencia habitual para expandir consultas sobre BM25. doc2query predice las preguntas que un documento puede responder y lo expande con ellas antes de indexarlo, lo que acerca la recuperación léxica a la de un reordenador neuronal con mucho menor costo en la consulta \citep{Nogueira2019Doc2query}. Con los modelos de lenguaje, HyDE genera un documento hipotético a partir de la consulta y busca los documentos reales cercanos a él \citep{Gao2023HyDE}, y Query2doc expande la consulta con un pseudodocumento generado, con mejoras de BM25 de entre el 3 y el 15\,\% \citep{Wang2023Query2doc}. Estas técnicas no ayudan siempre: en 24 modelos de recuperación y 12 colecciones, la ganancia de la expansión disminuye con la calidad del recuperador, y la expansión suele perjudicar a los modelos fuertes \citep{Weller2024}. En documentos financieros, HyDE aportó poco en las consultas numéricas precisas \citep{Akarsu2026}.

\textbf{Relación con JobAI.} JobAI no expande la consulta ni los documentos. La evaluación incluye una categoría de consultas de vocabulario, que mide el desajuste directamente, y la expansión queda como trabajo futuro, con la advertencia de \citet{Weller2024} de que el beneficio puede desaparecer sobre un pipeline que ya reordena.

\subsection{Reordenamiento con \textit{cross-encoders}}

Los modelos densos codifican la consulta y el documento por separado y los comparan con una similitud entre vectores; un \textit{cross-encoder} los lee juntos y estima su relevancia, con más detalle y mucho mayor costo. Por eso se aplica solo a los primeros candidatos de una etapa barata. \citet{Nogueira2019} mostraron que un reordenador basado en BERT superaba en un 27\,\% relativo en MRR@10 al mejor sistema previo de MS MARCO. En BEIR, los modelos de reordenamiento obtuvieron los mejores resultados promedio, a un alto costo computacional \citep{Thakur2021}. Los modelos de lenguaje generativos también ordenan: con la instrucción adecuada, ChatGPT y GPT-4 igualan o superan a los métodos supervisados, y esa capacidad puede destilarse en un modelo de 440 millones de parámetros \citep{Sun2023RankGPT}. Aun así, los \textit{cross-encoders} tradicionales siguen siendo muy competitivos frente a los reordenadores basados en GPT-4 \citep{Dejean2024}. Un reordenador tampoco es una mejora segura: con el mismo \textit{recall}, el nDCG@10 de un \textit{cross-encoder} pequeño cae hasta un 50\,\% cuando los distractores pasan de documentos al azar a negativos difíciles extraídos por el recuperador \citep{Lin2026Plug}. Hoy existen reordenadores multilingües de tamaño moderado, como bge-reranker-v2-m3, basado en BGE-M3 \citep{Chen2024M3}, y la serie Qwen3-Reranker, construida sobre modelos de lenguaje de la familia Qwen3 \citep{Zhang2025Qwen3}.

\textbf{Relación con JobAI.} El pipeline reordena las 20 primeras candidatas con Qwen3-Reranker-0.6B, elegido frente a bge-reranker-v2-m3 en un conjunto de desarrollo (Capítulo~\ref{cap:artefacto}); un reordenador generativo de gran tamaño habría multiplicado la latencia. La evaluación mide su aporte y su costo, y la dependencia de la primera etapa que señalan \citet{Lin2026Plug} es una razón para medirlo sobre las candidatas reales del pipeline y no de forma aislada.

\subsection{Modelos de \textit{embeddings} multilingües}

La familia E5 \citep{Wang2022E5} se entrena de forma contrastiva con pares de textos débilmente supervisados; según sus autores, fue el primer modelo que superó a BM25 en BEIR sin datos etiquetados. Su versión multilingüe \citep{Wang2024mE5} sigue la misma receta con mil millones de pares de textos en varios idiomas, en tamaños pequeño, base y grande. BGE-M3 \citep{Chen2024M3} cubre más de 100 idiomas y combina recuperación densa, dispersa y de múltiples vectores en un solo modelo. Estos modelos se comparan en MTEB, que abarca 8 tareas, 58 colecciones y 112 idiomas y en el que ningún modelo domina todas las tareas \citep{Muennighoff2023}, y en MIRACL, con más de 700.000 juicios de hablantes nativos en 18 idiomas, entre ellos el español, para la recuperación con la consulta y los documentos en el mismo idioma \citep{Zhang2023MIRACL}.

\textbf{Relación con JobAI.} Las ofertas del corpus están en inglés, pero los usuarios pueden escribir en español, un escenario entre idiomas que MIRACL no mide. La Iteración~2 reemplazó e5-base-v2, entrenado en inglés, por multilingual-e5-base, de la misma dimensión, y la evaluación compara ambos sobre las mismas consultas.

\section{Bases de datos vectoriales}
\label{sec:estado-arte-vectordb}

La recuperación densa requiere resolver la búsqueda de vecinos más cercanos sobre muchos vectores. \citet{Malkov2020} proponen HNSW, un grafo de proximidad organizado en capas jerárquicas con búsqueda de complejidad aproximadamente logarítmica, que hoy usan la mayoría de los motores vectoriales. \citet{Johnson2019}, con FAISS, aceleran en GPU la búsqueda exhaustiva y la aproximada mediante cuantización de producto. \citet{Aumuller2020}, con ANN-Benchmarks, proponen una metodología estandarizada para comparar estos algoritmos mediante curvas de compromiso entre \textit{recall} y velocidad.

Los algoritmos se ofrecen de dos formas. La primera es una extensión de un motor relacional, como pgvector sobre PostgreSQL, que agrega un tipo de dato vectorial e índices ANN. La segunda es un motor dedicado. \citet{Wang2021} describen con Milvus el diseño de esta segunda familia: separa almacenamiento y cómputo, admite varios tipos de índice y está pensada para cargas heterogéneas a gran escala. La revisión de \citet{Pan2024} organiza más de veinte de estos sistemas según su indexación, su soporte de búsqueda con filtros y su escalamiento, y la de \citet{Taipalus2024} señala como problemas abiertos la consistencia transaccional y la actualización incremental de los índices.

Esta literatura muestra que los motores dedicados están diseñados para escalas y patrones de acceso en los que una extensión relacional puede quedarse corta, en particular la búsqueda con filtros \citep{Sanca2024}. No establece, sin embargo, a partir de qué escala ni con qué selectividad la diferencia importa en la práctica, y una extensión relacional conserva ventajas operativas: transacciones, integridad referencial y una sola base de datos que administrar.

\textbf{Relación con JobAI.} JobAI almacenaba sus vectores en pgvector. La revisión del sistema mostró que su lentitud se debía a una consulta que no usaba el índice y no al motor (Capítulo~\ref{cap:artefacto}). Por eso este trabajo no da por justificada la migración a Qdrant, sino que la evalúa como parte del costo de la búsqueda: ambos motores quedan detrás de una interfaz común y se comparan con la metodología de \citet{Aumuller2020}, sobre datos reales y sintéticos, con y sin filtros (Sección~\ref{sec:met-preguntas}).

\section{Clasificación de consultas y text-to-SQL}
\label{sec:estado-arte-hibrido}

No todas las preguntas sobre un catálogo de ofertas son de recuperación. Una pregunta como \textit{\enquote{¿cuántas ofertas remotas de Python hay?}} tiene una respuesta exacta que da una agregación SQL; resolverla con una búsqueda vectorial, que devuelve los fragmentos más similares y no un conteo, produce una respuesta aproximada. \citet{Biswal2024} formalizan esta observación con TAG (\textit{Table-Augmented Generation}): el text-to-SQL puro y el RAG puro cubren cada uno solo una parte de las preguntas que un usuario hace sobre una base de datos, y ambos caminos deben combinarse.

La dificultad de generar SQL correcto con un modelo de lenguaje sin restricciones queda documentada por \citet{Li2023} con BIRD, un banco de 12.751 pares de pregunta y consulta sobre 95 bases de datos reales de 37 dominios. Incluso modelos muy capaces quedan muy por debajo del desempeño humano cuando deben generar SQL ejecutable sobre esquemas reales, con sus inconsistencias y su volumen. Esa brecha motiva clasificar la intención de la consulta antes de decidir si se consulta la base de datos, en lugar de delegar todo a un modelo de lenguaje.

\textbf{Relación con JobAI.} El clasificador de JobAI combina patrones léxicos con un clasificador basado en LLM que elige entre recuperación y agregación y, en el segundo caso, una plantilla de un catálogo cerrado de consultas SQL parametrizadas (Capítulo~\ref{cap:artefacto}). El modelo nunca genera SQL libre, lo que evita el riesgo que documenta \citet{Li2023}, a cambio de cubrir solo las preguntas previstas en el catálogo. La evaluación mide la exactitud de ese enrutamiento.

\section{Evaluación de la recuperación: colecciones de prueba y jueces}
\label{sec:estado-arte-juicios}

La evaluación \textit{offline} de la recuperación sigue el paradigma de Cranfield: un corpus fijo, un conjunto de necesidades de información y juicios de relevancia permiten comparar sistemas sin usuarios \citep{Sanderson2010}. Sus supuestos se han puesto a prueba durante décadas. Los evaluadores humanos discrepan en buena parte de los juicios, pero esas diferencias cambian poco el orden relativo de los sistemas, aunque sí sus cifras absolutas \citep{Voorhees2000}. Como no se puede juzgar todo el corpus, los juicios se construyen mediante agrupamiento: se juzga la unión de los primeros resultados de varios sistemas. Las medidas habituales no son robustas a juicios muy incompletos \citep{BuckleyVoorhees2004}, un agrupamiento poco profundo construido solo con sistemas tradicionales puede sesgar la evaluación contra los sistemas neuronales \citep{Yilmaz2020}, y una colección con demasiados documentos relevantes por tema deja de ser fiable \citep{Voorhees2022TooMany}. Con el mismo esfuerzo de juicio, muchos temas evaluados con poca profundidad dan más potencia estadística que pocos temas evaluados a fondo \citep{Webber2008}.

Juzgar a mano todos los pares agrupados es costoso, y la literatura reciente estudia el uso de modelos de lenguaje como jueces. \citet{Thomas2024} mostraron que un LLM con un prompt cuidadosamente diseñado puede predecir las preferencias de evaluadores humanos con una precisión comparable a la de estos, y \citet{Upadhyay2024} publicaron UMBRELA, una reproducción abierta de ese procedimiento. En la pista RAG de TREC 2024, los juicios de UMBRELA ordenaron los sistemas de forma muy correlacionada con los juicios humanos, aunque la correlación por tema fue menor \citep{Upadhyay2025}. Los LLM también pueden completar los huecos de un agrupamiento incompleto \citep{Upadhyay2024Holes}.

La evidencia en contra es igual de concreta. Los jueces LLM acuerdan más entre sí que con las personas y favorecen a los sistemas que usan LLM para ordenar: sobrestiman a los reordenadores basados en LLM entre 9 y 17 posiciones en promedio \citep{Frobe2025}. Pierden diferencias significativas que los juicios humanos detectan y producen falsas alarmas \citep{Sakai2025}, ordenan de forma injusta a los mejores sistemas y dan una tasa muy alta de falsos positivos en las diferencias estadísticas \citep{Otero2025}, y tienden a inflar la relevancia de pasajes que no satisfacen la necesidad \citep{Yu2026}. \citet{Faggioli2023} advierten sobre la circularidad cuando el juez y el sistema comparten modelo; \citet{Soboroff2025} sostiene que un juez LLM fija ese modelo como techo del desempeño medible, y \citet{ClarkeDietz2024} discuten que estos juicios puedan reemplazar a los humanos. La posición intermedia, que comparten varios de estos trabajos, es usar el juez LLM con una validación humana explícita y reportar sus límites.

\textbf{Relación con JobAI.} La evaluación usa un modelo de lenguaje como juez, con el procedimiento de UMBRELA adaptado al dominio, y lo valida contra juicios humanos anotados sin ver ningún juez, como hace LinkedIn con su juez de producción \citep{Borisyuk2026}. Los juicios se construyen por agrupamiento de profundidad 10 con todos los sistemas comparados, tradicionales y neuronales, y se juzgan todos los pares agrupados, de modo que no quedan huecos dentro del agrupamiento. El sesgo que documentan \citet{Frobe2025} es una amenaza concreta, porque el reordenador del pipeline es un modelo de lenguaje, y el trabajo tiene un solo anotador humano (Sección~\ref{sec:met-validez}).

\section{Evaluación de sistemas RAG}
\label{sec:estado-arte-evaluacion}

Evaluar un sistema RAG exige separar la recuperación de la generación. Para la generación, \citet{Es2024} proponen RAGAS, un conjunto de métricas sin respuesta de referencia. La \emph{fidelidad} descompone la respuesta en afirmaciones y verifica cada una contra el contexto recuperado; la \emph{relevancia de la respuesta} mide si la respuesta atiende a la pregunta formulada. ARES entrena jueces ligeros con datos sintéticos para la relevancia del contexto, la fidelidad y la relevancia de la respuesta, y corrige sus errores con unos cientos de anotaciones humanas \citep{SaadFalcon2024}. RAGChecker propone métricas de diagnóstico para cada módulo, con mejor correlación con los juicios humanos que las métricas anteriores \citep{Ru2024RAGChecker}, y eRAG mide la calidad de la recuperación por la utilidad de cada documento para el generador \citep{Salemi2024eRAG}. En aplicaciones de reclutamiento, RAGAS ya se usa como evaluación principal \citep{Tran2025}.

Otra parte de la literatura evalúa por componentes y por casos difíciles. BERGEN estandariza los experimentos de RAG porque la variedad de configuraciones impide atribuir los resultados a cada componente \citep{Rau2024BERGEN}. En documentos financieros, un pipeline de dos etapas, con búsqueda híbrida y reordenamiento, superó a todos los métodos de una etapa \citep{Akarsu2026}. Las métricas de recuperación pueden mejorar mientras las de la respuesta final se comportan de forma no monótona, por lo que la evaluación debería incluir el análisis de fallas por etapa y no solo la exactitud de la respuesta \citep{Muthyam2026}. Los modelos de lenguaje toleran cierto ruido en el contexto, pero fallan al rechazar una pregunta cuando ningún documento la responde \citep{Chen2024RGB}, y ninguna configuración de RAG resulta óptima a la vez para las preguntas con respuesta y para las que no la tienen \citep{Peng2025}.

La evaluación de un sistema RAG también debe considerar cómo falla. \citet{Barnett2024} reportan la experiencia de construir tres sistemas RAG en dominios distintos e identifican siete puntos de falla: la información no está en el corpus, el documento relevante no queda entre los primeros, no llega al contexto, el modelo no la extrae, la respuesta tiene un formato equivocado, una especificidad incorrecta o está incompleta. Concluyen que la validación de un sistema RAG solo es posible durante su operación y que su robustez se construye de forma evolutiva.

\textbf{Relación con JobAI.} La evaluación mide la fidelidad y la relevancia de las respuestas con RAGAS, y además mide salidas intermedias del pipeline: la lista de ofertas de cada componente y el cumplimiento de las restricciones. La compuerta de sin resultados responde al problema del rechazo que documentan \citet{Chen2024RGB} y \citet{Peng2025}, y se evalúa con consultas fuera de dominio. Las fallas encontradas durante las iteraciones se registran y se relacionan con los puntos de falla de \citet{Barnett2024}, sin plantear su detección como una pregunta de investigación.

\section{Latencia percibida en la búsqueda}
\label{sec:estado-arte-latencia}

La latencia es un costo porque los usuarios la perciben. En la búsqueda web, un retraso añadido menor que 500~ms no se distingue, y uno mayor que 1.000~ms se nota con mucha probabilidad, sobre todo entre los usuarios de un sistema rápido \citep{Arapakis2014,Bai2017}. En la búsqueda móvil la tolerancia es mayor, pero los retrasos de más de 7 a 10~s empeoran claramente la experiencia \citep{Arapakis2021}. Son estudios de buscadores web, en los que la respuesta es una lista de resultados y no un texto generado.

\textbf{Relación con JobAI.} Los umbrales de \citet{Arapakis2014} y \citet{Bai2017} sirven de referencia para el retraso que agregan los componentes del pipeline, con la advertencia de que provienen de otro tipo de interfaz (Sección~\ref{sec:met-enfoque}).
